\chapter{Building Your Own Reports}
\label{ch:reports}

The system does not lock you into a fixed set of printouts. Anything it records,
you can ask for in your own layout, save that layout, and run it again whenever
you like. Supervisors and administrators both have this screen, and it works the
same way for both.

\section{Building a report}

Open \menu{Reports}. You do not need to know anything technical --- you are
filling in a form.

\shot{22-admin-reports}{The report designer. Choose a dataset at the top, then
tick the columns you want.}{fig:rep-designer}

\begin{steps}
  \item Choose a \field{Dataset} --- the kind of information you want. For
        example \emph{Articles / WIP} for garments, or \emph{QC Inspections} for
        quality.
  \item Choose a \field{Report type}:
        \begin{bullets}
          \item \emph{Detail} gives one row per item.
          \item \emph{Summary} groups them and counts.
        \end{bullets}
  \item Tick the columns, or the groupings if you chose Summary.
  \item Add \field{Filters} to narrow it down --- for example
        \emph{Section is not SHIPPED}.
  \item Press \btn{Run report}.
\end{steps}

\shot{23-admin-reports-result}{A saved report, opened and run.}{fig:rep-run}

\begin{bullets}
  \item \btn{Export to CSV} saves the result for Excel.
  \item \btn{Save this report} keeps the layout so you or anyone else can run it
        again tomorrow.
\end{bullets}

\begin{tipbox}
Start from a saved report rather than a blank one. Press \btn{Open} on any report
in the list, change one thing, then \btn{Save as new}. It is far quicker than
building from scratch.
\end{tipbox}

\begin{plainwords}
A \textbf{dataset} is simply which list you are asking about --- garments,
inspections, handovers, fabric rolls. A \textbf{filter} is a condition that
throws away the rows you do not want. Nothing you tick here can damage anything:
a report only ever reads.
\end{plainwords}
